% ====================================================================
%  RESPONSE TO REVIEWER 1
%  Manuscript TP-1083 -- "The Incommensurability Principle in
%  Biological Transport" -- Transport Phenomena (De Gruyter Brill)
%  Standalone document; compile with: pdflatex Response-to-Reviewer-1
% ====================================================================
\documentclass[11pt]{article}
\usepackage[a4paper,margin=1in]{geometry}
\usepackage{amsmath,amssymb}
\usepackage[dvipsnames]{xcolor}
\usepackage{enumitem}
\usepackage{microtype}
\usepackage{hyperref}
\setlength{\emergencystretch}{2em}
\hypersetup{colorlinks=true,linkcolor=blue,urlcolor=blue}

% Reviewer 1 changes are shown in BLUE in the highlighted manuscript.
\newcommand{\rone}[1]{\textcolor{blue}{#1}}
\newenvironment{comment}%
  {\par\medskip\noindent\itshape\color{black!75} Reviewer comment.\par\nobreak}%
  {\par\nobreak\noindent\rule{\linewidth}{0.4pt}\medskip}
\newcommand{\response}{\par\medskip\noindent\textbf{Response.}\ }

\begin{document}

\begin{center}
{\large\bfseries Response to Reviewer 1}\\[4pt]
Manuscript \textbf{TP-1083} --- \emph{The Incommensurability Principle in
Biological Transport}\\
\emph{Transport Phenomena} (De Gruyter Brill)
\end{center}

\bigskip
\noindent We thank Reviewer~1 for the careful reading and for an observation that
improved the framing of the paper. Below we reproduce the comment (paraphrased
from the report) and detail our response. The highlighted copy uses a distinct
colour per source, as the editor requested: \rone{blue} for Reviewer~1, red for
Reviewer~2, green for editorial requests, and magenta for author-initiated
refinements.

\bigskip
\hrule
\bigskip

\noindent\textbf{Comment 1 --- Opening contrast (mouse vs.\ whale heart rate).}
\begin{comment}
The manuscript is well written and the concepts are well explained. The topic
is interesting and will attract Transport Phenomena readers. The introduction
contrasts the mouse ($\sim$600~bpm) and the whale ($\sim$20~bpm) heart rates.
Could the comparison be made in terms of the volume of blood pumped rather than
the heart rate?
\end{comment}

\response We agree. We have replaced the heart-rate contrast with a comparison
based on stroke volume and cardiac output. Cardiac output scales as
$\dot{Q}\propto M^{3/4}$ and stroke volume as $\propto M$, so a single whale
contraction ejects a blood volume roughly seven orders of magnitude larger than
the mouse's. That volume
travels through vessels whose radius $R\propto M^{3/8}$ sets the Womersley
number $\mathrm{Wo}\propto M^{1/4}$ (derived in Section~6)---making the whale's
vasculature the more inertial of the two despite its slower beat, which is
precisely the point. The cardiac-output allometry is now cited to Noujaim et
al.\ (\emph{Circulation} \textbf{110}, 2004, ``From Mouse to Whale'').
\emph{Location:} Section~1, \rone{blue}.

\end{document}
